\paragraph{One scenario, one domain.}
Every episode instantiates the same synthetic growth-agent situation with six
memories and five candidate actions. The construction was chosen so that the
corrupted belief is decision-relevant but not salient; other geometries ---
more memories, more symptoms, corrupted beliefs of intermediate relevance ---
are unexplored. Our claims are about the existence and direction of the
allocation bias in a controlled setting, not about its magnitude in any
production system.

\paragraph{Single-decision episodes.}
An episode is one decision with one round of verification, a deliberately
minimal probe of a long-horizon property. The lineage that produced the
drifted memory is simulated by construction rather than by running the
compressing agent for weeks. Whether intent-aligned allocation compounds
across sessions --- each session inheriting what the previous one failed to
check --- is the natural and harder follow-up.

\paragraph{The pilot's headline estimate moved under randomization.}
Our fixed-order pilot (45 episodes) showed one model verifying the corrupted
memory 5/5; under randomized order the same model's rate is 36/50. With
$n{=}5$ the pilot cannot statistically distinguish these (Fisher $p = 0.31$),
so we do not claim position confounded the pilot --- only that a fixed layout
leaves the question open, which is why the pre-registered experiment
randomizes presentation order and why we would caution against fixed-layout
memory-verification evaluations generally.

\paragraph{Harmful-decision rule is strict by design.}
The five-clause rule counts only unhedged, unverified, corrupted-evidence
rollouts as harmful. Models that hedge into guarded tests while retaining a
corrupted belief score as safe; the epistemic exposure we measure would then
surface later, outside the episode. Softer rules would report nonzero
behavioral failure; we chose the rule that gives models maximal credit and
committed to it in advance.

\paragraph{Verification is answered truthfully and instantly.}
The archive always returns the correct source record at zero latency and zero
noise. Real provenance is slower, sometimes missing, and sometimes itself
stale --- all of which would lower the value of verification and plausibly
sharpen the allocation problem.

\paragraph{Model coverage and settings.}
Six models from two providers, each at one reasoning setting, with structured
output enforced. Different effort settings, sampling temperatures, or
free-form output formats may shift allocation behavior; the between-model
heterogeneity we observe is itself a reason to expect sensitivity.
